This study was designed to address two related questions in heart failure prognosis: which clinical variables appear most strongly associated with mortality, and which predictive models perform most effectively in a small structured dataset. From the modeling perspective, the tree-based ensemble methods showed the strongest overall performance. RF produced the highest AUC, whereas XGBoost achieved the highest accuracy, precision, and F1-score in the main benchmark. Logistic regression remained competitive, while SVM, KNN, and the MLP baseline did not surpass the leading ensemble models. A plausible explanation is that the dataset is small, heterogeneous, and likely to contain nonlinear associations and threshold effects among laboratory and physiological variables. In such settings, ensemble tree methods are often well suited to capturing interactions without strong distributional assumptions. Their relatively stable behavior across preprocessing variants further supports their practical robustness in tabular clinical prediction tasks.

From the risk factor detection perspective, the multi-angle analysis produced a coherent pattern. Variables such as ejection fraction and serum creatinine were supported by both univariate statistical testing and model-based feature importance, which strengthens their interpretation as clinically meaningful markers of prognosis. This finding is medically plausible because ejection fraction reflects cardiac systolic function, whereas serum creatinine reflects renal impairment and the broader cardiorenal burden that is closely linked to adverse outcome in heart failure. The agreement between statistical comparison and model-derived ranking suggests that these variables capture a substantial portion of the mortality-related signal in the dataset.

At the same time, the prominent role of \texttt{time} requires careful interpretation. Follow-up time was highly informative in both the statistical and model-based analyses, but it is tied to the duration of observation and the structure of outcome ascertainment rather than functioning as a purely baseline biological determinant of risk. It should therefore be distinguished from clinically actionable baseline predictors when interpreting the combined ranking. Beyond hard class assignment, the framework also provides individualized mortality probabilities for new patients, which may offer a more interpretable summary of model output in practical use. More broadly, the study remains limited by its small sample size, reliance on a single benchmark dataset, and restricted variable set, all of which constrain generalizability and increase the risk of instability in both predictive estimates and factor ranking. The findings should therefore be interpreted as a structured benchmark and exploratory risk-factor assessment rather than as definitive evidence for clinical deployment. Future work should focus on larger multi-center cohorts, external validation, calibration analysis, and clearer separation between baseline predictors and follow-up-related variables.
